\section{Assumption discharge routes}
\label{app:dischargers}

\subsection{Stage complementarity and the derived tipping thresholds}
\label{sec:flip}

\emph{No result on the class under the lifting condition of
\Cref{sec:four-regimes} uses it; \Cref{prop:flip-route} below is what consumes it.}

As the political state hardens, both players' marginal payoffs from
sanctions can rise and eventually change sign. The following
complementarity condition turns this change into derived tipping
thresholds.

\begin{assumption}[Stage complementarity in $(s,x)$]
\label{ass:flip}
The stage payoffs are strictly supermodular in the sanction and the
political state, uniformly in the compliance action:
\begin{equation}
\frac{\partial^2 U_S}{\partial s\,\partial x}(x,a,s)>0
\qquad\text{and}\qquad
\frac{\partial^2 U_T}{\partial s\,\partial x}(x,a,s)>0
\qquad\text{for all }(x,a,s)\in[0,1]^3.
\label{eq:flip}
\end{equation}
\end{assumption}

Economically, a harder state raises the sender's marginal value of
sanctioning as sanction-supporting constituencies strengthen, and it
raises the target incumbent's marginal payoff from being sanctioned
through rally or habituation. The proofs use the corresponding
increasing-differences condition in derivative-free form \citep{topkis1998};
\eqref{eq:flip} is a smooth sufficient condition. It does not restrict
the curvature of the target payoff in sanctions or the sign of the
target's direct payoff from the political state. Nor does it, with the
other sign conditions, deliver persistence on its own: sign conditions
order slopes and bound no magnitude, and there is an economy satisfying
every one of them in which every equilibrium sanctions at zero at every
date (\Cref{app:formalization}); the magnitude
condition that persistence needs is the sub-pivot slope condition (P5) of
\Cref{thm:persistence}.

\paragraph{Flip sets and derived pivots.}
Fix a sanction-floor parameter $\bar s\in(0,1]$ and a positive stage
slope level $m_S$. On the general-$\F$ quantitative route, the
requirement that the stage gain outweigh the continuation change is
\[
\beta_S K_S\bigl|\F(x,s_2)-\F(x,s_1)\bigr|
<m_S(s_2-s_1)
\]
for the ordered pairs used in the basin proof. Under
\eqref{eq:affine}, this reduces exactly to
$\beta_S K_S\delta<m_S$. Denote the \emph{target flip set} by
$\mathcal F_T(\bar s)$ and the \emph{sender floor flip set at slope
level $m$} by $\mathcal F_S(\bar s;m)$:
\begin{align}
\mathcal F_T(\bar s)&:=\bigl\{x\in[0,1] :
\Ut(x,a,\cdot)\text{ is non-decreasing on }[\bar s,1]
\text{ for every }a\in[0,1]\bigr\},
\label{eq:FT-def}\\
\mathcal F_S(\bar s;m)&:=\bigl\{x\in[0,1] :
\Us(x,a,s_2)-\Us(x,a,s_1)\ge m(s_2-s_1)
\;\forall\,0\le s_1\le s_2\le\bar s,\;\forall a\in[0,1]\bigr\}.
\label{eq:FS-def}
\end{align}
Both sets quantify uniformly over compliance, so their definitions do
not depend on the equilibrium policy. The boundary sign conditions are
\begin{equation}
1\in\mathcal F_T(\bar s)\cap\mathcal F_S(\bar s;m_S),
\qquad
0\notin\mathcal F_T(\bar s)\cup\mathcal F_S(\bar s;m_S):
\label{eq:boundary}
\end{equation}
the flip has happened at the top state and not at the bottom.

\begin{proposition}[Derived tipping thresholds]
\label{prop:tipping}
Under \Cref{ass:flip} and \eqref{eq:boundary}, the pivots
\begin{equation}
\bar x_T:=\inf\mathcal F_T(\bar s),
\qquad
\bar x_S:=\inf\mathcal F_S(\bar s;m_S)
\end{equation}
exist, are interior ($0<\bar x_T,\bar x_S\le1$), satisfy their own
flip conditions, and the flip regions are exactly the upper
intervals
$\mathcal F_T(\bar s)=[\bar x_T,1]$,
$\mathcal F_S(\bar s;m_S)=[\bar x_S,1]$.
\end{proposition}

The basin is therefore derived rather than imposed: its component
thresholds are the infima of upper flip sets.

\Cref{prop:flip-route,app:dischargers} connect the derived basin to
\Cref{thm:persistence,thm:dominance,thm:shock}. The static and dynamic sender
thresholds, the per-action thresholds under the substitution between compliance and the sanction, an example
economy with interior pivots, and the pivots' comparative statics are in
\Cref{app:moved-ad}.

\subsection{Interval dominance and the derived-basin composition}

\begin{remark}[Interval dominance: the exact condition consumed]
\label{rem:ido}
A1 is the exact ordinal field consumed by the persistence certificate:
the interval dominance order of \citet{QuahStrulovici2009} for the
policy-substituted family
$\{\widetilde J_{\Vt,\Vs}(\cdot,x):x\in\basin\}$. It delivers a
greatest sender argmax that is non-decreasing in the state. Single
crossing in $(s,x)$ implies interval dominance, so every
single-crossing discharger remains sufficient. When the
continuation is non-monotone, interval dominance can instead be
obtained from single-peakedness in the sanction and an
escalation-consistency condition. The latter is automatic wherever
single crossing holds, but is a genuine primitive for general
value sets when the continuation slope is unrestricted. An example
economy discharges single-peakedness and escalation consistency even
though single crossing fails, establishing that the weaker route is
non-vacuous.
\end{remark}

\begin{proposition}[General-$\F$ route through the transition-adjusted slope condition]
\label{prop:F-chord-route}
Fix a Bellman-invariant set of value functions and pivots
$\bar x_T,\bar x_S$ with $\bar x=\max(\bar x_T,\bar x_S)<1$.
Suppose every sender value in the set is $K_S$-Lipschitz and
non-decreasing in the inherited state, the target combined slope condition holds,
and the canonical myopic ceiling holds. For every value pair in the set, let
$\alpha$ denote its induced target policy and assume that, for every
$x\in[\bar x_S,1]$ and $0\le s_1<s_2\le1$,
\[
\Us(x,\alpha(x),s_2)-\Us(x,\alpha(x),s_1)
>
\beta_S K_S\bigl|\F(x,s_2)-\F(x,s_1)\bigr|.
\tag{GF-Chord}\label{eq:GF-chord}
\]
Together with the row-contraction conditions for existence, these
hypotheses construct the local general-$\F$ invariant-set certificate and
hence deliver persistence, the dynamic-over-myopic comparison, and the
general minimum-drift shock result.
\end{proposition}

The exact-displacement condition does two jobs. On the low sanction range its
non-negative right-hand side implies strict stage-payoff monotonicity,
which pins the level of the dynamic argmax. On the basin, its
$K_S$-Lipschitz continuation bound makes the whole dynamic sender
objective strictly increasing, hence supplies the ordinal
interval-dominance field. No modulus of $\F$ in its sanction coordinate
is introduced. The separate hypothesis that the sender's values in the invariant set are monotone is
still load-bearing for the general-$\F$ headline certificate's
dynamic-over-myopic comparison; \eqref{eq:GF-chord} does not silently
remove it. If \textup{(F6-Q)} and the composite stage-state Lipschitz
bound hold, then $K_S$ in \eqref{eq:GF-chord} can be set to the closed
form \eqref{eq:KS-closed}.

The same exact-displacement argument survives the atomless moving-average
shock.  The box smear preserves the base value's $K_S$-Lipschitz
constant, so the shocked sender objective is strictly increasing when
the stage gain dominates
$\beta_S K_S\operatorname{dist}(\operatorname{clip}\F(x,s_1),
\operatorname{clip}\F(x,s_2))$.  This yields shocked single crossing
and monotonicity of the greatest shocked sender best response without
an affine transition or a separate sanction-coordinate modulus for
$\F$.  The exact-displacement condition is therefore the paper-facing
general-$\F$ shocked route.

\begin{proposition}[Derived-basin route to
\Cref{thm:persistence,thm:dominance,thm:shock}]
\label{prop:flip-route}
Let the transition follow the general law \eqref{eq:transition} and rise
under full sanctions below the derived basin, $x<\F(x,1)$ for
$x\in[0,\bar x]$ with $\bar x=\max(\bar x_T,\bar x_S)$ (so $\bar x<1$).
Suppose the sender half of \Cref{ass:flip} on the floor range $[0,\bar s]$
(its derivative-free form suffices), the sender's top boundary sign
$1\in\mathcal F_S(\bar s;m_S)$ in \eqref{eq:boundary} (the other signs are
needed only for the pivot interiority of \Cref{prop:tipping}, not here),
\Cref{ass:sender} on the derived basin, row-sum contraction of the joint
Bellman operator on a pair of invariant sets of value functions, $\beta_T+L_{TS}<1$ and
$\beta_S+L_{ST}<1$ (stronger than the Schur condition of \Cref{sec:existence}),
and also:
\begin{enumerate}[label=(\roman*)]
\item \emph{(above-floor diagonal)} at every state $x\ge\bar x_S$ and every
value pair in the invariant sets, the sender stage slope at the induced target policy clears the
margin on the above-floor range:
$\Us(x,\alpha(x),s_2)-\Us(x,\alpha(x),s_1)\ge m_S(s_2-s_1)$ for
$\bar s\le s_1\le s_2\le\max(x,\bar x)$;
\item \emph{(monotone sender values)} every sender value
$\Vs\in\mathcal C_S$ is non-decreasing in the state;
\item \emph{(myopic ceiling)} at every basin state the myopic equilibrium
sanction is at most $\max(x,\bar x)$.
\end{enumerate}
Then conclusions (i)--(iii) of \Cref{thm:persistence}, the statewise
ordering of the myopic sanction below the forward-looking one, and
the shock entry of \Cref{thm:shock} hold with the basin
$[\max(\bar x_T,\bar x_S),1]$ given by the derived pivots of
\Cref{prop:tipping}. No target preference or sign condition is used
(\Cref{ass:target}, the target half of \Cref{ass:flip}, the target's top
sign, monotone target values in the invariant set): the target enters only through its
operator's row bound and its policy in the invariant sets, and its flip set only as
the location $\bar x_T$ of the basin's lower end.
\end{proposition}

The three goals read neither the target half of \Cref{ass:flip}, nor the target's
top sign, nor monotone target values in the invariant set, nor \Cref{ass:concavity}. With the target half,
the target's top sign and monotone target values in the invariant set, the target combined
slope condition of \Cref{ass:target} is still a derived result: the stage slope is
non-negative above the target pivot and the continuation slope is
non-negative by monotonicity of those values; it is what conclusion (vi) of
\Cref{thm:persistence} would need. The residual clauses, the general-$\F$ route of
\Cref{prop:F-chord-route}, the discharger for a positive stage slope, and the route's
scope limits are detailed
below. The route is inhabited by an example economy with pivots
$\bar x_T=2/5$ and $\bar x_S=1/2$, and the example is a family over the
transition: its conclusions hold at every transition satisfying
\textup{(F1)--(F4)} whose full-sanction section is $\F(y,1)=y/2+1/2$ and whose
lower bound on the sanction slope is at least $11/36$
(\path{sfFCert_three_goals_generalF}). The affine law at rate $\tfrac12$ is one
member (\path{sfCert_three_goals_generalF}); the quadratic-bump transition, with
lower bound $47/100$ on the sanction slope and affine at no rate, is another
(\path{sfqbCert_three_goals_generalF}, \path{sfqbParams_F_ne_affineF}). The
corner section is an equality on the line $s=1$, so at every member the
realized full-sanction path is the affine one. These are existence
certificates: they show that the route does not need the affine law,
and they say nothing about transitions outside the family. The example's kinked
sender payoff is
essential because the diagonal-residual and myopic-ceiling clauses are
incompatible near the top state for a smooth sender payoff with a
Lipschitz derivative.

\paragraph{Dischargers for the residual clauses.}
Clause (iii) of \Cref{prop:flip-route} has a stage-level discharger, and
so has clause (i) when $\bar x_T\le\bar x_S$. The discharger of (i) gives the
diagonal on $[\bar s,x]$ at states $x\ge\bar x_S$ above $\bar s$, the form the
earlier affine route read; this is clause (i) when $\bar x_T\le\bar x_S$,
since then $\max(x,\bar x)=x$, while for $\bar x_T>\bar x_S$ the slope on
$[\bar s,\bar x]$ at states in $[\bar x_S,\bar x)$ is not derived. That
diagonal is derived from
$\alpha$-antitonicity on the basin (available under the
substituted-objective slope condition discussed in
\Cref{sec:mechanism}), the action-complementarity sign
$\partial^2\Us/\partial s\partial a\le0$, stage complementarity, and
a single \emph{pivot-action frontier} condition: margin-level slope conditions
at the one action $\alpha(\bar x_S)$ along the diagonal states, so
the per-basin-state quantification collapses to a one-action family.
The myopic ceiling (iii) follows from strict stage decrease above the
floor at the myopic action (``the myopic action is on the un-flipped
side of the diagonal''); the bridge from this stage-level condition
to fully primitive form (the $\sigma$-difference between the stage
and dynamic best responses) remains open. Complementing these, the
next-state reparametrization $y=\F(x,s)$, under which the
continuation cancels from every slope comparison, gives monotonicity
of the greatest next-state selection per fixed action from stage
complementarity plus concavity of $\Us$ in $s$ alone (the Amir
1996 pattern), and under strict complementarity \emph{every} selection from the stage
argmax is monotone in the state.

\paragraph{Deriving the dynamic single crossing: the route-(b)
convexity cone.} The one sender-side dynamic input the
primitives of the derived-basin route provably cannot deliver for free is the
single crossing of the \emph{forward-looking} objective
$\Us(x,\alpha(x),s)+\beta_S\Vs(\F(x,s))$ in (sanction; state):
single crossing is not preserved under sums, and at the affine
transition the continuation term has the requisite increasing
differences precisely when $\Vs$ is convex along the transition
range (Topkis Thm 2.6.2). It is \emph{derivable} from a
range-local curvature cone: if every sender value in the value set is
convex along the basin transition range $[\,\mathrm{lo},1]$ with
$\mathrm{lo}\le(1-\delta)\bar x$, then together with stage
complementarity, the compliance-placates sign, and
$\alpha$-antitonicity, the dynamic objective has increasing
differences (hence single crossing, and a fortiori the certificate's
interval-dominance condition of \Cref{rem:ido}).
The cone's economic content is \emph{no diminishing returns to
entrenchment on the post-tipping basin}: each further step of
political hardening is worth weakly more than the last, the
increasing-returns lock-in mechanism of the institutional
path-dependence literature \citep{Pierson2000}, imposed only where
that literature claims it. The contested and low-basin regions are
unconstrained, so dual-basin path dependence is untouched. Two caveats:
the cone excludes saturation of entrenchment returns strictly
inside the basin (if saturation is the preferred semantics, the
derivative-free form with its three existing dischargers remains canonical);
and preservation of the cone under the Bellman map cannot remain
range-local (invariance of convexity on $[\,\mathrm{lo},1]$ pulls in
the strictly larger interval $[(1-\delta)\mathrm{lo},1]$ and
iteration expands the requirement to all of $[0,1]$), so a
\emph{constructive} convex-cone invariant set would require the global
commitment rejected above; the form stated here, which places the cone on a chosen value set, is
the endpoint of the range-local route.
Constant-$\Vs$ value sets (e.g.\ that of the auxiliary polynomial economy of
\Cref{tab:witness-roster}, whose sender value is constant) lie in the cone
with equality.

\paragraph{The Lipschitz route under a positive stage slope: single crossing
\emph{constructed} at the equilibrium.} Because the convexity cone of
the previous route is near-vacuous against the \emph{concave} value
functions standard dynamic programming delivers
\citep[Thm~4.8]{StokeyLucas1989}, a second
route discharges the single crossing for the genuinely \emph{non-monotone}
equilibrium value, using only Lipschitz regularity. Suppose the sender
stage payoff has a strictly positive lower slope $m_S>0$ in $s$ on the
\emph{full} choice axis $[0,1]$ at the equilibrium compliance (the
\emph{positive-stage-slope} setting, in which the marginal sanction benefit stays
positive across $[0,1]$), and that the transition-adjusted slope condition
\eqref{eq:GF-chord} holds, with $K_S$ the Lipschitz constant of the invariant set.
Then the dynamic objective is strictly increasing in $s$ on $[0,1]$: the
stage rise strictly exceeds the worst-case continuation fall
$\beta_S K_S|\F(x,s_2)-\F(x,s_1)|$, an inequality that uses
\emph{only} that $\Vs$ is $K_S$-Lipschitz, never that it is monotone or
convex. Under the affine law, the lower slope $m_S$ and
$\beta_S K_S\delta<m_S$ are a sufficient specialization. Strict
monotonicity in $s$ is the corner property (the sender plays the
high-sanction boundary), and it implies the single-crossing clause of
\Cref{ass:sender} a fortiori.

This is a property the \emph{constructed} equilibrium
\emph{has}, not one imposed on it. The invariant set is the closed
$K_S$-Lipschitz ball; the Bellman operator maps it into itself and the
ball is closed in sup-norm, so the contraction's fixed point inherits
$K_S$-Lipschitzness by construction
\citep[Cor.~1 to Thm~3.2]{StokeyLucas1989}. Single crossing is therefore
\emph{derived} from the positive-stage-slope condition and \eqref{eq:GF-chord},
both stated over the model's existing primitives, and the certificate's
interval-dominance field (\Cref{rem:ido}) is filled by this derivation
(the derived single crossing upcasts for free) rather than supplied as
a separate hypothesis, so every high-basin result consuming the
certificate inherits a derived condition. Unlike the convexity cone,
this route admits the non-monotone equilibrium value the model actually
produces.

The route's scope limit travels with it. It is confined to the
positive-stage-slope setting; in the complementary
setting (where the sender's stage payoff \emph{falls} in $s$ and
only the forward-looking continuation gain carries the high-sanction
corner), the two-sided $\pm K_S$ bound is too coarse to capture that
gain at a non-monotone $\Vs$, so in that setting the
interval-dominance field of \Cref{rem:ido} remains assumption-backed
rather than derived (\Cref{sec:flip}). The derivation above should not
be read apart from this limit.

\begin{remark}[Lipschitz strict-slope margin as a cardinal sufficient
condition]
\label{rem:wedge-equivalence}
A cardinal sufficient condition for \Cref{ass:sender} is the
Lipschitz strict-slope margin
\[
U_S(x,a,s_2)-U_S(x,a,s_1)
>
\beta_S K_S\lvert\F(x,s_2)-\F(x,s_1)\rvert
\]
for $s_1<s_2$ in $[\max(x,\bar x),1]$. Combined with
$\Vs\in\mathcal C_S$, this
strict difference-quotient inequality implies the Milgrom-Shannon single-crossing
clause of \Cref{ass:sender} (the stage gain strictly dominates the
worst-case Lipschitz absorption of the continuation, so the increment
of $J_S$ in $s$ is bounded below by $\mathrm{gap}\cdot(s_2-s_1)>0$;
hence single-crossing is satisfied with the constant-sign upper
witness). The cardinal form is thus available to users with explicit
numerical bounds, and it is more than an optional convenience: the
\emph{level} conclusion of \Cref{lem:single-crossing}, that no
sender maximizer lies below $\max(x,\bar x)$, is proved from this
same strict difference-quotient inequality imposed on the \emph{complementary}
range $[0,\max(x,\bar x)]$ (equivalently, from the margin on the whole
of $[0,1]$), which bounds the stage gain against the worst-case
Lipschitz absorption of a possibly non-monotone $\Vs$. Under the
affine law, a lower slope $m_S$ and
$\beta_S K_S\delta<m_S$ imply this exact inequality. The ordinal
form is the canonical statement on which the proof otherwise depends.
\end{remark}

\begin{remark}[Fixed-action interval dominance does not survive policy
substitution automatically]
\label{rem:policy-substitution-obstruction}
Continuity and feasibility of the target policy are not enough to turn
fixed-action interval dominance into \Cref{ass:sender}. The
machine-checked counterexample takes
\[
  G(s,x,a)=(x-a)s,\qquad \alpha(x)=2x,
\]
on sanctions \([0,1]\) and states \([0,\tfrac12]\). For every fixed
action \(a\), the family \(G(\cdot,\cdot,a)\) satisfies interval
dominance, and \(\alpha\) is continuous and feasible. After
substitution, however,
\[
  G(s,x,\alpha(x))=-xs,
\]
which violates interval dominance. A valid discharge must therefore
derive an appropriate order or variation restriction on the realized
policy, or verify interval dominance directly after substitution. The
convex-value route remains valid precisely because it derives policy
antitonicity before proving the realized objective's single crossing.
\end{remark}
